\section{Pulse Length}
\label{sec:pulse-length}

The pulse duration imposes an additional constraint on spectroscopic resolution, as the spectral width of the drive in the frequency domain scales inversely with the pulse length. Figure 8 presents a numerical calculation of the extracted linewidth obtained using a square pulse, showing that for short pulse durations (  $\leq 20~\mu s$), the measured spectroscopic linewidth is dominantly limited by Fourier broadening rather than decoherence.

In contrast, for pulse durations exceeding approximately 20 $\mu s$, the echo-based pulse sequence approaches the decoherence-limited regime, yielding spectroscopy with an $\mathcal{O}(1)$ agreement with the $T_2$-limited linewidth.

For longer pulse durations, an additional effect emerges: the off-resonant background signal gradually decays while the resonance feature itself remains approximately constant. This results in the formation of a characteristic narrow peak nested within the broader echo-induced dip, reflecting the onset of adiabatic state evolution in the vicinity of the resonance.

